\documentclass[../../main.tex]{subfiles}

\begin{document}

\section{集合的分离与函数延拓}\label{section:ch21-分离延拓}

在前一节我们证明了局部紧 Hausdorff 空间的几个基本性质. 我们现在证明这些性质对于在紧 Hausdorff 空间上的 Borel 集上构造测度是有用的. 然而, 存在对于这些应用所需要的分离与延拓的某些性质, 我们现在陈述它们为一个定理. 首先定义以下两个术语.

\begin{definition}[Borel 集的内外正则性]\label{definition:ch21-内外正则}
令 $\mu$ 为紧 Hausdorff 空间 $X$ 上的 Borel 测度, $E$ 是 Borel 集.
\begin{enumerate}[(i)]
\item 集合 $E$ 称为是关于 $\mu$ 外正则的, 若
\begin{align*}
\mu(E)=\inf\{\mu(\mathcal{O})\mid E\subseteq\mathcal{O},\ \mathcal{O}\text{ 是开的}\}.
\end{align*}
\item 集合 $E$ 称为是关于 $\mu$ 内正则的, 若
\begin{align*}
\mu(E)=\sup\{\mu(K)\mid K\subseteq E,\ K\text{ 是紧的}\}.
\end{align*}
\item 集合 $E$ 称为是关于 $\mu$ 正则的, 若它既是外正则的又是内正则的.
\end{enumerate}
\end{definition}

\begin{definition}[Radon 测度]\label{definition:ch21-radon}
紧 Hausdorff 空间 $X$ 上的 Borel 测度 $\mu$ 称为是 Radon 测度, 若每个 Borel 集是关于 $\mu$ 正则的.
\end{definition}

\begin{theorem}[分离与延拓定理]\label{theorem:ch21-分离延拓}
令 $X$ 为紧 Hausdorff 空间, 其 Borel $\sigma$ 代数 $\mathcal{B}(X)$ 是 Baire $\sigma$ 代数 $\mathcal{B}_{0}(X)$ 的扩张. 则 $X$ 满足以下性质:
\begin{enumerate}[(i)]
\item \textbf{紧集的分离.} 若 $K$ 是 $X$ 的紧子集而 $\mathcal{O}$ 是包含 $K$ 的开集, 则存在集合 $K_{1}$ 和 $\mathcal{O}_{1}$ 使得
\begin{align*}
K\subseteq\mathcal{O}_{1}\subseteq K_{1}\subseteq\mathcal{O},
\end{align*}
其中 $\mathcal{O}_{1}$ 是开 $F_{\sigma}$ 集, 而 $K_{1}$ 是闭 $G_{\delta}$ 集.
\item \textbf{Baire 集的延拓.} 若 $K_{0}$ 是 Baire 集, 则存在 Borel 集 $K$ 使得只要 $\mu$ 是 $\mathcal{B}(X)$ 上的 Borel 测度且 $\nu$ 是 $\mathcal{B}_{0}(X)$ 上的 Radon 测度使得 $\nu=\mu|_{\mathcal{B}_{0}(X)}$, 就有 $\mu(K)=\nu(K_{0})$.
\end{enumerate}
\end{theorem}

\subsection*{习题}

\begin{exercise}\label{exercise:ch21-22}
证明: 若 $X$ 是紧度量空间, 则 $\mathcal{B}(X)=\mathcal{B}_{0}(X)$.
\end{exercise}

\begin{exercise}\label{exercise:ch21-23}
证明: 若 $X$ 是紧 Hausdorff 空间, 则每个开集是紧集的可数并的闭包.
\end{exercise}

\begin{exercise}\label{exercise:ch21-24}
令 $X$ 为紧 Hausdorff 空间. 证明 Borel $\sigma$ 代数 $\mathcal{B}(X)$ 等于 Baire $\sigma$ 代数 $\mathcal{B}_{0}(X)$ 当且仅当 $X$ 是可数多个紧 $G_{\delta}$ 集的并.
\end{exercise}

\end{document}